\documentclass[11pt]{article}
\usepackage[margin=1in]{geometry}
\usepackage{longtable}
\usepackage{array}
\usepackage{multirow}
\usepackage{hyperref}
\usepackage[T1]{fontenc}

\setlength{\parindent}{0pt}
\setlength{\parskip}{8pt}

\begin{document}

\begin{center}
{\large \textbf{Stat 32960: Topics in dimensionality reduction, Winter 2027}} \\
\end{center}




\textbf{Course description} \quad 
This course surveys modern statistical methods and theory for estimating low-dimensional structure in high-dimensional datasets. Tentative topics include high-dimensional theory of PCA, structured matrix factorization (sparse PCA, NMF, Bayesian methods), nonlinear techniques and manifold learning (principal curves, ISOMAP, tSNE, UMAP), autoencoders, and approaches for multi-view data and tensors. Emphasis is placed on the practical application of these methods to real-world high-dimensional datasets.

\medskip

\textbf{Prerequisites} \quad Graduate student in Statistics or related field, strong background in linear algebra, and a
strong interest in research. 

\medskip

\textbf{Course info}

\begin{itemize}
  \item Instructor: Phillip Nicol, \texttt{[pnicol@uchicago.edu]} (Office hours: TBD)
  \item GitHub repository containing all materials: 
\end{itemize}

\textbf{Assignments}

\begin{itemize}
  \item There are 3 homework assignments due throughout the quarter. Some HW exercises are open-ended and at the level of current research problems in the field. For these, complete solutions are not required, but they should be attempted.
  \begin{itemize}
      \item Late HW is accepted, provided an extension is arranged before the due date.
  \end{itemize}
  \item There is 1 in-class paper presentation.
  \item The final project involves extending an idea discussed in class to a new scenario. There are different options depending on the student's interests:
  \begin{itemize}
      \item Theory: Extend an existing result to a new setting, or analyze a model in a new theoretical regime. 
      \item Methods: Design or extend a dimension reduction technique to address a new and relevant challenge. For example, this could entail considering a new data modality or a scientific question. 
  \end{itemize}
  The deliverable will be a presentation (likely during the final exam period) as well as a brief writeup. Students may work individually or in groups of up to three. 
  
\end{itemize}


\textbf{Grading} \quad The final grade is based on the following:
\begin{itemize}
  \item 25\% Participation, 25\% Homework, 25 \% Paper presentation, 25 \% Final project
\end{itemize}

\textbf{Collaboration and AI policy}

\begin{itemize}
  \item Students are strongly encouraged to discuss homework problems with classmates, but each student should turn in their own homework assignment. 
  \item Students may use AI tools on homework assignments, provided that the final writeup is their own work. In particular, AI generated text or math is not permitted in the submitted homework. AI generated computer code (e.g., to run simulations) is permitted. 
  \item AI use is permitted and encouraged on the final project. However, the final writeup cannot contain AI generated math or text. 
\end{itemize}

\newpage

\textbf{Schedule (tentative)}

\medskip

\begin{itemize}
  \item \textbf{Week 1} Introduction, PCA 
  \item \textbf{Week 2} PCA (Cont.), Perturbation theory
  \item \textbf{Week 3} Spiked covariance theory
  \item \textbf{Week 4} Projection pursuit
  \item \textbf{Week 5} Structured matrix factorization
  \item \textbf{Week 6} Multidimensional scaling, ISOMAP
  \item \textbf{Week 7} tSNE/UMAP
  \item \textbf{Week 8} Autoencoders
  \item \textbf{Week 9} Multi-view data, tensors
\end{itemize}

\end{document}
